\begin{center}\begin{tikzpicture}[font=\small,>=stealth,line join=round]
\draw[gray!40,rounded corners=3pt] (0,0) rectangle (14.2,8.3);\node at (7.1,7.700000000000001) {证明三：从任意起点出发，同一过程处处成立};
\begin{scope}[shift={(1.7153846153846155,3.8842857142857143)},x=0.7692307692307693cm,y=1.3142857142857143cm]\draw[->,gray] (-0.8,0)--(7,0) node[right]{$x$};\draw[->,gray] (0,-1.7)--(0,1.8) node[above]{$y$};\draw[blue,thick,solid,domain=-0.4:6.6,samples=180] plot(\x,{sin(deg(pi*\x/2))});\fill[orange!85!black] (0.5,0.7071067811865476) circle(2.2pt);\fill[orange!85!black] (2.5,-0.7071067811865476) circle(2.2pt);\fill[orange!85!black] (4.5,0.7071067811865476) circle(2.2pt);\draw[orange!85!black,dashed] (.5,.70710678)--(4.5,.70710678);\node[above=5pt] at (.5,.70710678){$P$};\node[below=5pt] at (2.5,-.70710678){$Q$};\node[above=5pt] at (4.5,.70710678){$R$};\end{scope}\node[anchor=west,align=left,text width=5.0cm] at (8.4,4.0) {图取 $x=\frac12$\\[9pt]$P=(x,y)$\\[9pt]$Q=(x+2,-y)$\\[9pt]$R=(x+4,y)$\\[9pt]$P,R$ 横向相距4};
\node[text width=12.6cm,align=center] at (7.1,.65) {示例中任意改变起点，前后两次反号的关系仍然成立。};
\end{tikzpicture}\end{center}
